%\documentclass{article}
%\usepackage[utf8]{inputenc}
%\usepackage{amsmath}

%\begin{document}

%\noindent \underline{\textbf{Problema 2.17:}}
%\vspace{0.3cm}

Consideremos um pêndulo balançando num fluido viscoso. Além das 6 variáveis usuais do pêndulo ($\ell, g, m, T_0, \theta$ e $T$), vamos considerar a viscosidade do fluido, $\mu$, a densidade de massa do fluido, $\rho$, e o diâmetro do pêndulo, $d$, cujas dimensões são
$$[\mu] = \frac{M}{LT}, \quad [\rho] = \frac{M}{L^3} \quad \text{e} \quad [d] = L.$$

Assim, o número total de variáveis é $n = 9$ e o número de dimensões primárias é $m = 3$ ($L, T$ e $M$). Então, pelo Teorema Pi de Buckingham, precisamos de $n - m = 6$ grupos adimensionais para relacioná-las. Tomemos $\ell, g$, e $m$ como as três variáveis fundamentais contendo $L, T$ e $M$. Daí, os grupos são
$$
\begin{aligned}
\Pi_1 &= \ell^{a_1} g^{b_1} m^{c_1} T_0; & \Pi_4 &= \ell^{a_4} g^{b_4} m^{c_4} \mu; \\
\Pi_2 &= \ell^{a_2} g^{b_2} m^{c_2} \theta; & \Pi_5 &= \ell^{a_5} g^{b_5} m^{c_5} \rho; \\
\Pi_3 &= \ell^{a_3} g^{b_3} m^{c_3} T; & \Pi_6 &= \ell^{a_6} g^{b_6} m^{c_6} d.
\end{aligned}
$$

Em termos das dimensões primárias, obtemos:

\vspace{0.5cm}
\noindent \textbf{i)} $\Pi_1 = L^{a_1} \left(\frac{L}{T^2}\right)^{b_1} M^{c_1} T = L^{a_1+b_1} T^{1-2b_1} M^{c_1} \\ \\ \Rightarrow \begin{cases} a_1 + b_1 = 0 \Rightarrow a_1 = -b_1 = -1/2 \\ 1 - 2b_1 = 0 \Rightarrow b_1 = 1/2 \\ c_1 = 0. \end{cases}$

Logo,
$$\Pi_1 = \ell^{-1/2} g^{1/2} m^0 T_0 = \sqrt{\frac{g}{\ell}} T_0 \Leftrightarrow \Pi_1 = \frac{T_0}{\sqrt{\ell/g}}.$$

\vspace{0.5cm}
\noindent \textbf{ii)} $\Pi_2 = L^{a_2} \left(\frac{L}{T^2}\right)^{b_2} M^{c_2} 1 = L^{a_2+b_2} T^{-2b_2} M^{c_2} \\ \\ \Rightarrow \begin{cases} a_2 + b_2 = 0 \Rightarrow a_2 = -b_2 = 0 \\ -2b_2 = 0 \Rightarrow b_2 = 0 \\ c_2 = 0. \end{cases}$

Logo,
$$\Pi_2 = \ell^0 g^0 m^0 \theta \Leftrightarrow \Pi_2 = \theta.$$

\vspace{0.5cm}
\noindent \textbf{iii)} $\Pi_3 = L^{a_3} \left(\frac{L}{T^2}\right)^{b_3} M^{c_3} \frac{ML}{T^2} = L^{a_3+b_3+1} T^{-2b_3-2} M^{c_3+1} \\ \\ \Rightarrow \begin{cases} a_3 + b_3 + 1 = 0 \Rightarrow a_3 = 0 \\ -2b_3 - 2 = 0 \Rightarrow b_3 = -1 \\ c_3 + 1 = 0 \Rightarrow c_3 = -1. \end{cases}$

Logo,
$$\Pi_3 = \ell^0 g^{-1} m^{-1} T \Leftrightarrow \Pi_3 = \frac{T}{mg}.$$

\vspace{0.5cm}
\noindent \textbf{iv)} $\Pi_4 = L^{a_4} \left(\frac{L}{T^2}\right)^{b_4} M^{c_4} \frac{M}{LT} = L^{a_4+b_4-1} T^{-2b_4-1} M^{c_4+1} \\\\ \Rightarrow \begin{cases} a_4 + b_4 - 1 = 0 \Rightarrow a_4 = 3/2 \\ -2b_4 - 1 = 0 \Rightarrow b_4 = -1/2 \\ c_4 + 1 = 0 \Rightarrow c_4 = -1. \end{cases}$

Logo,
$$\Pi_4 = \ell^{3/2} g^{-1/2} m^{-1} \mu \Leftrightarrow \Pi_4 = \frac{\mu}{m} \sqrt{\frac{\ell^3}{g}}.$$

\vspace{0.5cm}
\noindent \textbf{v)} $\Pi_5 = L^{a_5} \left(\frac{L}{T^2}\right)^{b_5} M^{c_5} \frac{M}{L^3} = L^{a_5+b_5-3} T^{-2b_5} M^{c_5+1} \\ \\ \Rightarrow \begin{cases} a_5 + b_5 - 3 = 0 \Rightarrow a_5 = 3 \\ -2b_5 = 0 \Rightarrow b_5 = 0 \\ c_5 + 1 = 0 \Rightarrow c_5 = -1. \end{cases}$

Logo,
$$\Pi_5 = \ell^3 g^0 m^{-1} \rho \Leftrightarrow \Pi_5 = \frac{\rho \ell^3}{m}.$$

\vspace{0.5cm}
\noindent \textbf{vi)} $\Pi_6 = L^{a_6} \left(\frac{L}{T^2}\right)^{b_6} M^{c_6} L = L^{a_6+b_6+1} T^{-2b_6} M^{c_6} \\ \\ \Rightarrow \begin{cases} a_6 + b_6 + 1 = 0 \Rightarrow a_6 = -1 \\ -2b_6 = 0 \Rightarrow b_6 = 0 \\ c_6 = 0. \end{cases}$

Logo,
$$\Pi_6 = \ell^{-1} g^0 m^0 d \Leftrightarrow \Pi_6 = \frac{d}{\ell}.$$

%\end{document}